% ============================================================================
\section{Introduction}
\label{sec:intro}
% ============================================================================

\subsection{Background and problem}
The College of Computing and Information Sciences (CoCIS) keeps valuable equipment that must be
protected from unauthorised access, and it pays high electricity bills, partly because lights are
left on in empty rooms. In October 2026 management agreed to pilot an automated system on the
\textbf{main entrance} and \textbf{two lecture rooms}.

\subsection{What we built}
A complete, self-contained control and monitoring system on the ESP32-S3, simulated in Wokwi and
written in C with ESP-IDF~v6.0 and FreeRTOS:
\begin{itemize}
  \item \textbf{Entrance terminal.} Two MFRC522 RFID readers, one at the entry and one at the
        exit, each with its own servo door lock; anti-passback; an opening-hours schedule that also
        holds people inside after hours; an LCD and a buzzer; and two-factor admin login (the admin
        card at a reader + a PIN).
  \item \textbf{Two lecture-room terminals.} A PIR motion sensor and a light sensor per room drive a
        dimmable lamp. Lights turn off when the room is empty or when there is enough daylight, and
        the room reports the energy it used and saved.
  \item \textbf{Storage and management.} Every event is written to the SD card as a
        \emph{hash-chained} (tamper-evident) log line, and a serial command menu shows who is
        inside, room status, lighting and the access log and lets an authenticated admin change
        the schedule, users and overrides.
\end{itemize}
The system runs entirely offline. An \emph{optional} web dashboard (Next.js + PostgreSQL) was also
prepared as an extension (\cref{sec:innovation}); the system does not depend on it.

\subsection{How this report maps to the required deliverables}
\begin{center}
\begin{tabularx}{\linewidth}{lL}
\toprule
Deliverable (brief \S4) & Where \\
\midrule
1. System operation report & \Cref{sec:requirements} to \cref{sec:limitations}: architecture,
   hardware, software design, operation of every feature, synchronisation, testing \\
2. Wokwi simulation link & Title page and \cref{sec:members} \\
3. Individual video demonstrations & \Cref{sec:videos} \\
4. Individual contributions & \Cref{sec:contributions} \\
5. Group member details & Title page and \cref{sec:members} \\
\bottomrule
\end{tabularx}
\end{center}

\subsection{Conventions}
Requirements from the brief are written \req{2.1} \ldots\ \req{2.7} (the brief's section numbers).
Paths such as \file{components/core_logic/access_policy.c} are relative to \file{firmware/} unless
they start with \file{dashboard/} or \file{docs/}. Default timings and thresholds quoted in this
report are the values in \file{main/Kconfig.projbuild}; all of them can be changed with
\code{idf.py menuconfig} without editing code.
